\section{Prediction log}\label{app:prereg}
Each job's predictions are in the header of its script on the public \texttt{gpu} branch, committed before the machine
started. Each machine cloned the published branch at boot and ran the header as committed, so git history orders every
prediction before its measurement. Of the \rtRentals{} rentals of jobs 093--111, every commit preceded its rental by
at least \rtCommitMin\,s, and every machine started its job at least \rtMinLag\,s after the push. For \rtNear{} of
them the push and the rental's creation fall within a second of each other, and for \rtAfter{} GitHub records the push
up to \rtAfterMax\,s after the rental was created (\texttt{scripts/reg\_timing.py}). \Cref{tab:prereg} summarises
every job's clauses, and \cref{tab:failures} gives the failures. Outcome notes with the statistics are in
\texttt{prereg/}; the batched-verification variants of \cref{sec:horizon} (rejected drafts free, routed like the true
token, or routed at random) are in \texttt{prereg/batchk\_outcome.md}. Jobs before 073, which built and debugged the
engine, also registered predictions in their headers. Their outcome notes are in \texttt{prereg/}, and they are not
scored here; job 072's, for one, failed or did not apply on two of its three.

\paragraph{Scoring.} Every prediction of jobs 073--111 is split into its separable clauses, and each clause is scored
under one rule: \emph{held} when its point estimate is on the predicted side and the 95\% paired-bootstrap interval
excludes the threshold (for a band, lies inside it); \emph{held (point)} when the point estimate is on the predicted
side but no interval exists or the interval crosses the threshold; \emph{failed} when the point estimate is on the
wrong side; \emph{untested} when the clause could not be evaluated; and \emph{void} when its job was voided. Jobs
073--\scLastJob{} were scored by hand from the outcome notes, into \texttt{prereg/scorecard\_clauses.json}, which
generates the scorecard. The hand scoring makes one exception: deterministic counts such as reads have no sampling
error, so it holds them outright, not held (point). From job 099 on, a script scores each job's clauses, without that
exception.

\paragraph{Two tallies.} The hand-scored tally covers the \scClauses{} clauses of jobs 073--\scLastJob: \scHeld{}
held, \scHeldPoint{} held on the point estimate, \scFailed{} failed, \scUntested{} were untested and \scVoid{} void.
By era, \scEarlyHeld{} of \scEarlyScored{} scored clauses held in jobs 073--075, when the system was being built;
\scMidHeld{} of \scMidScored{} in jobs 076--081, \scMidSign{} of them sign-only predictions made after a same-machine
pilot; \scLateHeld{} of \scLateScored{} in jobs 082--087; and \scRevHeld{} of \scRevScored{} in jobs 088--\scLastJob.
The last row of \cref{tab:prereg} is the overall tally, for jobs 073--111. The hand-scored tally is a subset of it:
the table's rows for jobs 073--\scLastJob{} are the hand scores, and the script-scored jobs 099--111 make up the rest.
In the overall row, a deterministic count is therefore held in jobs 073--\scLastJob{} and held (point) after.

\paragraph{Scoring checks.} A script (\texttt{scripts/scorecard\_auto.py}) re-scores every hand-scored clause whose
measured value is a number and whose threshold parses, \scaAuto{} of \scaN. It agrees with the hand score on
\scaAgree. Of the \scaDisagree{} differences, \scaDet{} are the deterministic counts (reads, divergences, table
entries) that the hand scoring holds outright and the script scores held (point), and \scaOther{} is a boundary value
the hand scoring read unrounded (1.249 against a band starting at 1.25). Of the clauses the script scores held on the
point estimate, \scaNoInterval{} have no interval and \scaCrosses{} have one that crosses the threshold. The
re-scoring is a check; neither tally uses it.

% generated by scripts/tab_prereg.py
\begin{table}[p]\centering\scriptsize
\caption{Every registered job from 073 on and how its clauses scored: \emph{held} with a 95\% interval on the predicted side (or outright, for deterministic counts), held on the point estimate only, failed, and untested or void. Jobs 073--098 were scored by hand, the rest by script; the scorecard lists every clause. \emph{Cond.}: a job's validity condition (at least $k$ valid machines), met or not; it is operational, not a prediction, and is not counted in the other columns.}\label{tab:prereg}
\setlength\tabcolsep{3pt}\begin{tabular}{@{}llrrrrc@{}}\toprule
Job & What it tested & Held & Point & Failed & Untested & Cond. \\\midrule
073 & ours against FreeToken, same host & 3 & 0 & 4 & 0 & -- \\
074 & the 40\% budget fixed & 0 & 2 & 4 & 0 & -- \\
074b & server overhead & 2 & 1 & 0 & 1 & -- \\
075 & RTX PRO 6000 at 40--60\% & 2 & 0 & 1 & 0 & -- \\
076 & the system table on an RTX 5090 & 2 & 1 & 0 & 0 & -- \\
077 & competitors, gpt-oss & 3 & 0 & 0 & 0 & -- \\
078 & Qwen3, four systems & 3 & 0 & 0 & 0 & -- \\
079 & Qwen3 profile & 1 & 3 & 0 & 0 & -- \\
080 & the fetch split & 4 & 0 & 0 & 0 & -- \\
081 & the system table with the calibrated fetch split & 6 & 0 & 0 & 0 & -- \\
082 & held-out prompts and steady state (review run 1) & 4 & 0 & 3 & 1 & -- \\
083 & gpt-oss, long outputs on hard held-out problems & 5 & 0 & 0 & 0 & -- \\
084 & all-in-VRAM column and traces (voided: throttled card) & 0 & 0 & 0 & 4 & -- \\
084b & all-in-VRAM column and traces, rerun & 2 & 2 & 1 & 0 & -- \\
085 & a slower PCIe link, gpt-oss & 5 & 0 & 1 & 0 & -- \\
085b & a slower PCIe link, Qwen3 & 5 & 0 & 1 & 0 & -- \\
086 & Qwen3.6 & 1 & 1 & 2 & 0 & -- \\
087 & ablation with a static cache & 2 & 1 & 2 & 0 & -- \\
088 & slow link, fair comparison & 5 & 3 & 3 & 0 & -- \\
089 & the system table on a second host & 9 & 7 & 12 & 0 & -- \\
090 & output parity against all-in-VRAM (RTX PRO 6000) & 10 & 0 & 0 & 0 & -- \\
091 & the grid on an RTX 4090 & 11 & 3 & 5 & 0 & -- \\
092 & the grid on an RTX 3090 & 6 & 1 & 7 & 1 & -- \\
093 & first foresight oracles & 9 & 2 & 21 & 4 & -- \\
094 & oracles that minimise bytes & 12 & 1 & 24 & 1 & -- \\
095 & single-read oracles & 44 & 0 & 12 & 0 & -- \\
096 & the oracle factorial & 29 & 97 & 12 & 0 & -- \\
097 & the learned admission order & 11 & 53 & 20 & 0 & -- \\
098 & FreeToken tuned & 22 & 15 & 7 & 0 & -- \\
099 & the host panel & 190 & 42 & 39 & 2 & -- \\
100 & windows of foresight & 41 & 122 & 27 & 26 & -- \\
101 & the crossover by link ratio & 6 & 19 & 0 & 0 & -- \\
102 & the read-schedule microbenchmark & 0 & 28 & 5 & 0 & -- \\
103 & fewest-admission schedule (plan not run) & 0 & 4 & 0 & 41 & -- \\
104 & fewest-admission schedule & 1 & 22 & 24 & 0 & -- \\
105 & \cref{eq:sum}, plain form; layer-ahead copy; fewest-admission set & 7 & 20 & 4 & 0 & -- \\
106 & \cref{eq:sum} with overlap; admitting less; slow links by round & 21 & 112 & 48 & 0 & -- \\
107 & \cref{eq:sum} on new machines; admitting less; fewest-admission loads & 11 & 22 & 12 & 0 & not met \\
108 & \cref{eq:sum} on new desktop-class machines & 0 & 12 & 6 & 2 & met \\
109 & the 2$\times$2 and policies without foresight on new machines & 61 & 29 & 0 & 0 & met \\
110 & the RTX 5090 trend on RTX 4090s (stopped by V1) & 0 & 0 & 0 & 1 & not met \\
111 & job 110 relaunched, V1 corrected & 24 & 13 & 0 & 0 & met \\
112 & higher-ratio RTX 4090s; timing control; second probe & 24 & 21 & 12 & 0 & met \\
113 & the two-read path with the in-step fetches off & 6 & 6 & 2 & 0 & met \\
114 & the 2$\times$2 with the in-step fetches off & 17 & 14 & 2 & 0 & met \\
115 & the same on more machines; idle probe runs & 11 & 16 & 1 & 0 & met \\
\midrule all & & 638 & 695 & 324 & 84 & 7 met, 2 not \\
\bottomrule\end{tabular}\end{table}


\paragraph{Failed predictions.} Most bands that failed in jobs 073--\scLastJob{} are the sizes of effects whose
direction held. In jobs 088--\scLastJob, \scRevFailBand{} of \scRevBandClauses{} bands failed, as did \scRevFailThr{}
threshold, \scRevFailEq{} equality and \scRevFailSign{} sign clauses; \scRevSignHeld{} of \scRevSignClauses{} sign
clauses held with an interval, and \scRevSignHeldAny{} held at least on the point estimate. \Cref{tab:failures} gives
the failures job by job: those that matter for jobs 073--\scLastJob{} (the supplement lists all \scFailed), and the
families of failures of jobs 099--111.

% The failed predictions of the prediction log (app_prereg.tex), job by job. Measured values are macros generated by
% the scoring scripts; the Predicted column holds the registered bands, thresholds and signs.
\begin{footnotesize}
\setlength\tabcolsep{3pt}\setlength\LTcapwidth{\textwidth}
\newcommand\tfjob[2]{\addlinespace[3pt]\textbf{#1} & \multicolumn{4}{>{\raggedright\arraybackslash}p{\dimexpr0.835\textwidth+6\tabcolsep\relax}@{}}{#2}\\}
\begin{longtable}{@{}>{\raggedright\arraybackslash}p{0.075\textwidth}>{\raggedright\arraybackslash}p{0.25\textwidth}>{\raggedright\arraybackslash}p{0.15\textwidth}>{\raggedright\arraybackslash}p{0.17\textwidth}>{\raggedright\arraybackslash}p{0.265\textwidth}@{}}
\caption{Failed predictions, by job. \emph{Predicted}: the registered band, threshold or sign. Jobs 073--\scLastJob{}
were scored by hand, and the table gives the failures that matter; from job 099 on, each job's line gives its tally as
its script scored it. ``---'': the reason, where known, is in the job's outcome note. The supplement lists every
clause.}\label{tab:failures}\\
\toprule
Job & What was predicted & Predicted & Measured & Why it failed, or what it means \\\midrule
\endfirsthead
\toprule
Job & What was predicted & Predicted & Measured & Why it failed, or what it means \\\midrule
\endhead
\midrule\multicolumn{5}{r}{\emph{continued}}\\
\endfoot
\bottomrule
\endlastfoot
\multicolumn{5}{@{}l}{\emph{Scored by hand, jobs 073--\scLastJob}}\\
\tfjob{085b}{A slower PCIe link, Qwen3.}
& Our lead over FreeToken larger than in \cref{tab:headline}, at 2 or more of 3 budgets (the job's own restatement of
its fourth prediction) & at least 2 of 3 & 1 of 3 & Downgraded from the running log; the pooled 4 of 6 held \\
\tfjob{088}{A slow link, fair comparison.}
& Our lead over FreeToken & 1.10--1.50$\times$ & \slFtMin--\slFtMax$\times$ & Above the band; the lead itself held \\
\tfjob{089}{The system table (\cref{tab:headline}) rerun on a second machine.}
& Our ratio over FreeToken, against host B's & within $\pm$0.06 & \scRatioFailMin--\scRatioFailMax{} below, at
\scRatioFailCells{} of six cells & --- \\
& Our speed at the host-bound cells, against host B's & within $\pm$12\% & \scOursHybDropMin--\scOursHybDropMax\%
lower & --- \\
& LRU slower than decayed frequency, at the two GPU-heavy cells & 10--30\% slower & \scLruFailMin--\scLruFailMax\%
slower & --- \\
& LRU still ahead of FreeToken at gpt-oss 11\% & ahead & behind & A sign failure \\
\tfjob{091, 092}{The grid on an RTX 4090 and on an RTX 3090 (the supplement).}
& Our lead over FreeToken; our share of the limit & 10--35\% (gpt-oss 11\%), 0--15\% (Qwen3); 25--45\% & above the
bands & FreeToken's carried-over backend fell far below its band \\
& Ours at least 1.8$\times$ llama.cpp, on the slow machines & at least 1.8$\times$ & 1.56--1.78$\times$ & --- \\
& The 4090 machine's FETCH tables & the registered entries & one more expert copied, in two tables & An equality
failure \\
& Ours at least 1.3$\times$ llama.cpp on Mixtral & at least 1.3$\times$ & a tie & --- \\
& The two launch orders agree & within 3\% & 7\%; a second difference, 11\% at gpt-oss 11\% on the 3090 machine, was reported, not scored & --- \\
\tfjob{093}{The first oracle job (\cref{sec:foresight}).}
& The hit-optimal oracle faster than the online policy, by 20--80\% & faster; 20--80\% & slower at gpt-oss 11\% and
Qwen3 12.5\%, the two lowest budgets & A sign failure as well as a band failure \\
& Its speed at the four host-bound cells & 45--65\% of the limit & 43.7--44.8\% & Just below the band \\
& Its hit rate and admissions, against the optimum's & within 3 points; within 30\% & outside, at three and at two
of the four cells & --- \\
& The cost of no overlap & 5--25\% & 0\% & --- \\
& The cost of all-CPU, Qwen3 & 0--10\% & 15\% & --- \\
& The hit-optimal oracle's gain at the two GPU-bound cells & below 15\% & 48 and 56\% & --- \\
\tfjob{094}{Oracles that minimise bytes: the bands missed as \cref{app:foresight} describes.}
& The bytes-optimal oracle beats the hit-optimal one, at the two lowest budgets & by at least 15\% & by less & --- \\
& Its 16-token window keeps its gain & at least 70\% of it & less, at three cells & --- \\
\tfjob{095}{Single-read oracles: 44 of 56 clauses held, and the 12 that failed are of size.}
& The model on \MinOne{}'s counters & within 10\% & 13--17\% off, at four cells & Size \\
& \MinOne{}'s reads, against the simulation's & within 10\% & 17\% below, at the two GPU-bound cells & The
registered simulator gave a resident carried over from the previous sequence no next use, and so evicted it first;
corrected, the simulation is within \fsfSimReadsDiffMax\% of the engine's reads at every cell (\cref{app:foresight})
\\
& The read-ahead oracle at the host-bound cells & 55--85\% of the limit & 44 and 51\% at the two 25\% cells & Size \\
& The model's over-prediction of the read-ahead oracle & 15--50\% & 7 and 13\% at the two lowest budgets & Size \\
& The paced variant, against the unpaced one & within 10\% & 12--13\% faster at the two 25\% cells & Size \\
\tfjob{096}{The oracle factorial on O4 and O5: the clauses held at most cells of both machines, and
\cref{app:foresight} lists the failures.}
& \MinOne{}'s gain, O4 against O5 & differs by at most 0.15 & by more, at five cells & The largest failure of the
job: the difference between the machines \\
& Belady's single-read prefetch, Qwen3 12.5\% on O4 & at most 1.05$\times$ the online policy & above & --- \\
\tfjob{098}{FreeToken tuned (\texttt{prereg/freetoken\_outcome\_098.md}): its best setting stayed within the predicted
10\% of its default, and our ratios within the predicted bands.}
& FreeToken's best backend, gpt-oss 25\% & hybrid & offload & An equality failure \\
& Stock llama.cpp against FreeToken's CPU backend, CPU only (gpt-oss; Qwen3) & at least 2$\times$; at least
1.2$\times$ & about equal & --- \\
& The fixed fetch caps; the thread count & within 8\% of the calibrated split; 8 threads within 3\% of the default or
slower & outside & As the job's outcome note describes \\
\addlinespace[3pt]
\multicolumn{5}{@{}l}{\emph{Scored by script, jobs 099--113}}\\
\tfjob{099}{The host panel (\texttt{scripts/panel\_099.py}; \texttt{prereg/scorecard\_099.json},
\texttt{prereg/panel\_outcome\_099.md}): \scbN{} clauses; \scbHeld{} held, \scbPoint{} held (point), \scbFailed{}
failed, \scbUntested{} untested.}
& Bands set from four earlier machines, none of them link-starved & the registered bands & most failures, on the two
machines whose link reads at less than half their CPU's rate & The bands did not allow for dependence on the
machine \\
& The planning time of the noisy whole-future window & at most 150\,$\mu$s per step & above & --- \\
\tfjob{100}{Windows of foresight (\texttt{scripts/job100.py}; \texttt{prereg/window\_outcome\_100.md}): \sccN{}
clauses; \sccHeld{} held, \sccPoint{} held (point), \sccFailed{} failed, \sccUntested{} untested.}
& The deployed-path windows miss within 10\% as often as a replay that admits at once & within 10\% & more misses;
most of the job's failures & Their copies land late \\
& The probe-only model, on the half-width-link machine & within 12\% & outside & --- \\
& Four other bands or orderings & as registered & at their edges & --- \\
\tfjob{101}{The crossover by link ratio (\texttt{scripts/job101.py}; \texttt{prereg/crossover\_outcome\_101.md}):
\jbClauses{} clauses; \jbHeld{} held, \jbPoint{} held (point), \jbFailed{} failed, \jbUntested{} untested.}
\tfjob{102}{The read-schedule microbenchmark on three machines (\texttt{scripts/job102.py};
\texttt{prereg/readsched\_outcome\_102.md}): \jdClauses{} clauses; \jdHeld{} held, \jdPoint{} held (point; none has
an interval), \jdFailed{} failed, \jdUntested{} untested.}
& Waiting per layer costs time against the analytic schedule & at least 0.03 of the host term & nothing measurable
on two of the three machines & Four of the five failures \\
& The order of the two budgets, Intel machine & as registered & the other way & --- \\
\tfjob{103}{MIN's fewest-admission schedule against the greedy one; the plan did not run (\texttt{scripts/job103.py}):
\jeClauses{} clauses; \jePoint{} held (point; the relaunch of Pf), \jeUntested{} untested.}
\tfjob{104}{Job 103 repeated, with the plan run (\texttt{prereg/minadm\_outcome\_104.md}): \jfClauses{} clauses;
\jfHeld{} held, \jfPoint{} held (point), \jfFailed{} failed.}
& \FewOne{}'s copies and misses at gpt-oss 25\%, against \MinOne{}'s & copies at most 0.70$\times$; misses within 3\%
& \jfCopyRatioMax$\times$; \jfMissDevMax\% more & --- \\
& The engine's misses, against the machine's replay & within 2\% & more & --- \\
& \FewTwo{}'s reads, against \MinTwo{}'s & at most 0.90$\times$ & more & --- \\
& The 2$\times$2's interaction, with the plan's set & shrinks & grew & --- \\
& \FewTwo{} against \MinTwo{} & at most 0.01 behind & 0.047 behind, on Pg at gpt-oss 25\% & Held at the other 5 of 6
cells \\
\tfjob{105}{The plain form of \cref{eq:sum} with $G$ profiled on each machine, \LA{} and \FewOne{}
(\texttt{scripts/job105.py}; \texttt{prereg/sumlaw\_outcome\_105.md}): \jgClauses{} clauses; \jgHeld{} held,
\jgPoint{} held (point), \jgFailed{} failed. Every failure is the relation at gpt-oss 25\%.}
& The relation at gpt-oss 25\%, per machine & within 6\% & over-predicts the deployed cache by
\jgErrMidFailMin--\jgErrMidFailMax\% on \jgErrMidFailN{} of four machines (\jgErrMidHeldMin\% on Pf) & At 11\% it is
within \jgErrLowMin{} to \jgErrLowMax\% \\
& Its pooled median error & at most 4\% & \jgErrMed\% & The 25\% cells take it past the band \\
\tfjob{106}{The relation of \cref{eq:sum} with its overlap term on both models at two budgets, the two online
admission rules, and MIN's sets launch by launch, on five machines (\texttt{scripts/job106.py}; \texttt{prereg/onlineadmit\_outcome\_106.md}):
\jiClauses{} clauses; \jiHeld{} held, \jiPoint{} held (point), \jiFailed{} failed: \jiFailedUnst{} on the two unstable
machines, \jiFailedStable{} on the \jiStable{} stable ones and \jiFailedPooled{} pooled.}
& Clauses on the two server machines new to the job & as registered & \jiFailedUnst{} failed & One computed wrong
outputs; the other varied by more than the registered 2\% between rounds (\cref{app:more}) \\
& \Margin{}'s speed-up at gpt-oss & at least 1.01 & below, on two stable machines & --- \\
& The learned order's speed-up & at least 1.00 & below, on the fast-memory machine & --- \\
& The learned order's host time & at most 400\,$\mu$s (gpt-oss), 700\,$\mu$s (Qwen3) per step & above, on the
fastest-link machine & --- \\
& The relation's prediction of both rules' ratios, Pf at 25\% & within 0.03 & off by 0.030 and \jiRatioPredMax & --- \\
& The relation's median error over all machines & at most 4\% & above & --- \\
\tfjob{107}{The relation, \Margin{} and MIN's fewest-admission set loaded both ways, on machines never rented before
(\texttt{scripts/job107.py}; \texttt{prereg/newhosts\_outcome\_107.md}): on its \jlValidWord{} valid machines,
\jlClauses{} clauses; \jlHeld{} held, \jlPoint{} held (point), \jlFailed{} failed.}
& The relation, per cell & within 6\% & beyond, at \jlFailCells{} of \jlCells{} cells & --- \\
& The relation's pooled median error; the overlap form against the plain one & at most 4\%; overlap form better &
above; not better & --- \\
& \Margin{}'s speed-up & at least 0.995 & below, at \jlDkFailWord{} cell & --- \\
& The relation's prediction of \Margin{}'s speed, on the server & within 0.03 & outside, at \jlDkPredFailN{} cells &
--- \\
& Valid machines & at least \jlRequiredWord & \jlValidWord & Inconclusive by its registered rule \\
\tfjob{108}{The relation on machines never rented before, gpt-oss only (\texttt{scripts/job108.py};
\texttt{prereg/smallhosts\_outcome\_108.md}): \jmClauses{} clauses; \jmPoint{} held (point), \jmFailed{} failed, and
the rest could not be tested.}
& The relation on the EPYC 7543, both budgets & within 8\% & outside & --- \\
& Its implied read rate, both budgets & 0.85--1.25 of $B_{\mathrm{host}}$ & outside & --- \\
& Pooled: the cells within 6\%; the median error & at most one cell beyond 6\%; at most 4\% & \jmWithinSix{} of
\jmCells{} within 6\%; \jmErrMed\% & --- \\
\tfjob{109}{The 2$\times$2 of \cref{sec:gap}, the read-ahead oracle and four variants without foresight, on desktop-class
machines never rented before (\texttt{scripts/job109.py}; \cref{tab:job109}): \olClauses{} clauses;
\olClausesHeldCI{} held with a 95\% interval over problems, \olClausesPoint{} held (point), \olClausesFailed{}
failed. Of its \olPredN{} predictions \olPredHeld{} held; one clause of P8, for ratios of at least 0.9, had no machine
to test it.}
\tfjob{110}{The RTX 5090 trend of \cref{sec:secondcard}, frozen in its header, on RTX 4090 machines: void as
registered.}
& Valid machines & at least 3 & none & The output check used the RTX 5090's reference loss \\
\tfjob{111}{Job 110 relaunched with the RTX 4090's own reference: job 110's clauses on \cxNWord{} valid machines;
\cxClausesHeld{} held with an interval, \cxClausesPoint{} held (point), \cxClausesFailed{} failed. Its sixth
prediction, for machines with a ratio of at least 0.5, had none to test.}
\tfjob{112}{The second card at higher ratios, the timing control, a second probe and two relaunches of job 109's
machines (\texttt{scripts/job112.py}; \cref{tab:timing,tab:job109}): \tcClauses{} clauses; \tcClausesHeld{} held with
an interval, \tcClausesPoint{} held (point), \tcClausesFailed{} failed.}
& \FewEarly's misses & at most 0.95 of \FewTwo's & \tcMissRatioLowRng{} at 11\%, \tcMissRatioMidRng{} at 25\%; failed on \tcTOneFailed{} of \tcTOneN{} machine-budgets & The replay of the publication rule predicted 0.91; the engine's fetch table keeps MIN's set from forming \\
& \FewEarly{} no slower than \FewTwo & at least $-0.01$ & \tcLagLowRng{} at 11\%, \tcLagMidRng{} at 25\%; failed on \tcTTwoFailed{} of \tcTTwoN & The early copies admitted and read more; the control cannot separate landing from reading \\
\tfjob{113}{The two-read path with the in-step fetches off (\texttt{scripts/job113.py}; \cref{tab:job113}): \hsClauses{}
clauses; \hsClausesHeld{} held with an interval, \hsClausesPoint{} held (point), \hsClausesFailed{} failed.}
& The deployed cache with the fetches off & 0.80 to 1.00 of the deployed cache & \hsBaseZeroLowRng{} at 11\%; failed on the \hsHThreeFailedOn & The probe-set fetch table slows the deployed cache on that machine \\
& One read over the fetch-free two-read arm & at least 0.05 & \hsOneReadMarginLowRng{} at 11\%; failed on the \hsHSixFailedOn & Once the set is held, the second read costs little there \\
\end{longtable}
\end{footnotesize}


\paragraph{Changes against the running log.} Against the running log kept during the study, two clauses are
downgraded. A llama.cpp clause of job 074b that the log counted as held was never run, and job 085b's own restatement
of its fourth prediction, ``at least 2 of 3'', failed at 1 of 3, though the pooled 4 of 6 held. Job 073's ``identical
outputs'' is scored on fully identical outputs (0--3\%) rather than on agreement of the first 160 characters
(27--47\%). Ten clauses the log counted as held, among them job 087's first, hold on the point estimate only under the
interval rule. The file \texttt{prereg/scorecard\_outcome.md} lists every change.

\paragraph{Reruns and rentals without results.} Job 084 ran on a throttled card and was rerun, with its predictions
unchanged, as 084b/084c. Job 085's Qwen3 half ran as 085b on a second machine with a slow link after a hung
conversion. Rentals that produced no results were destroyed and are in the rental ledger
(\texttt{gpu/vast\_ledger.json}): two of job 097's four rentals; the first two for 099a and the first for 099d,
replaced before any measurement; 100d, which stalled downloading the model; 100e, which never reached its job; 103c,
which never started its job; and 105c, 105d and 105g below.

\paragraph{Jobs 099--104.} Job 099's two untested clauses are one machine's replay check, whose step failed when a
Python package did not install. Job 100's header gained a ninth prediction, about the two replacement machines, after
the first four machines had started; it compares the replacements with values one of which was already measured. Job
100's untested clauses are the planning time of the deployed-path windows, which the engine does not time, and the
paired prediction for a machine that never reached a timed step. Job 101 relaunched the lowest-ratio panel machine and
added a second Ryzen 7 9800X3D behind a slower link. In job 103 the machines could not install the plan's solver
beside Debian's numpy, so the engine ran the greedy schedule in the plan configurations, as it logs. Job 104 repeated
job 103 with the plan in its own environment on three of the same machines: Pf; a second Core Ultra 9 285K, whose link
reads at 0.49 of its CPU rate; and Pg. Job 103's 5950X had turned out to be Pg's second launch. Most of job 104's
speed predictions held: \FewOne{} gains on Pf, and gains at least 0.03 over \MinOne{} at gpt-oss 11\% on every
machine.

\paragraph{Jobs 105 and 106.} Two of job 105's clauses predicted quantities already measured: the profiled $G$ (job
069c) and the plan's copies, which are a property of the trace (job 104). They test that the engine reproduces them.
Every speed prediction of job 105 held: \LA{} gains at least 3\% on the machine with a ratio above 0.8 and loses on
both below 0.4, and \FewOne{} beats the deployed cache on every machine, beats \MinOne{} by at least 0.03 on the three
with a ratio under 0.7, and copies at most 0.65 of \MinOne{}'s experts. Three more launched machines produced no
results: 105c never started its job, and the job's bandwidth gate stopped 105d and 105g, whose cards were throttled,
before any timed run. Job 106 ran on five machines. Two server machines were new to the job; one computed wrong
outputs, and the other varied by more than the registered 2\% between rounds (\cref{app:more}). The \jiStable{} stable
machines are Pf, the Threadripper 9960X of job 105 and panel machine Pe, all rented before, and on them the relation
held on every cell. \Margin{}'s pooled median speed-up of at least 1.02 held only because of the two unstable
machines; on the stable ones it is \jiDkPooled.

\paragraph{Jobs 107 and 108.} Job 107 had validity gates registered before launch. Of \jlLaunched{} machines rented,
\jlGated{} stopped at the gate and \jlInvalid{} failed the round check. The test had \jlValidWord{} valid machines of
the \jlRequiredWord{} it required, which by its registered rule makes it inconclusive; its per-cell clauses failed on
both. The profiled GPU compute, \Margin{}'s reads and pooled median, both fewest-admission loads and the crossover
between them held. Replacing machines that failed the round check, and the rules that chose the last two machines,
were amendments committed before each machine started. Jobs 102, 105 and 106 also had their headers edited after
launch, each time to substitute a machine, with the predictions unchanged and marked so in the header. Job 108 tested
the relation on machines never rented before, with the class drawn after job 107 gated before any download, gpt-oss
only, and an 8\% band. Of \jmLaunched{} machines rented, \jmGated{} stopped at the gate, \jmInvalid{} failed the round
check, and \jmValidWord{} were valid. The profiled $G$ of the EPYC 7543 could not be tested.

\paragraph{Job 109.} The job fixed its population, machines never rented before with a link-to-CPU ratio of at least
0.5, before any started; \olStartedWord{} machines started, and all \olNWord{} passed every gate and round check. A
smoke run on gpt-oss-20b, on a machine outside the sample, first checked that the two new combinations run and compute
the same model. One listed offer was no longer offered and was skipped as registered. The job's deadline cut the 25\%
round of the Ryzen 9 5950X, the slowest machine. The header's statement that no listed offer was in the rental ledger
was wrong for one: the Core i9-13900KF's offer had been rented twice before (the first rental for 099d, and 100e),
both times without a result. Its GPU had never been measured, and it passed the identity gate.

\paragraph{Jobs 110 and 111.} Job 110 tested the RTX 5090 trend of \cref{sec:secondcard}, frozen in its header, on RTX
4090 machines; \cxTenStarted{} machines started. Of these, \cxTenGated{} stopped at its ratio gate (at
\cxTenGatedRatio, below 0.25), and every other one failed its output check after one round. That check used the RTX
5090's reference loss. The RTX 4090 computes a loss \cxLossDiff\% higher in every configuration of every process, so
the check failed by design rather than by fault, and job 110 is void as registered. Job 111 relaunched three of the
four machines that had run a round, with the card's own reference (its first round's loss on one of them) and job
110's predictions unchanged. The fourth, 110e, was not relaunched: its GPU was left out of the relaunch list, an
oversight we found after the job. The predictions were committed before any RTX 4090 data, and the corrected reference
after one round on these machines. Of job 111's first launches, \cxElevenNoModel{} failed to download the model,
before any model run; they were relaunched once, by an amendment committed before the relaunch.

The clause-by-clause tables (\scClauses{} rows for jobs 073--098, \scbN{} for job 099, \sccN{} for job 100,
\jbClauses{} for job 101, \jdClauses{} for job 102, \jeClauses{} and \jfClauses{} for jobs 103 and 104, \jgClauses{}
for job 105, \jiClauses{} for job 106, \jlClauses{} for job 107, \jmClauses{} for job 108, \olClauses{} for job 109,
\cxClauses{} for jobs 110 and 111) are in the supplement (\texttt{paper/supplement.pdf}), generated by the same
scripts.
